\documentclass[10pt]{article}
\usepackage[margin=1in]{geometry}
\usepackage{amsmath,amssymb,amsthm,booktabs,hyperref}
\hypersetup{hidelinks}
\newtheorem{theorem}{Theorem}[section]
\newtheorem{lemma}[theorem]{Lemma}
\newtheorem{remark}[theorem]{Remark}
\newcommand{\Q}{\mathbf Q}
\newcommand{\F}{\mathbf F}
\title{Five-Term Arithmetic Progressions and Cubes in a Pure Cubic Field}
\author{\texttt{AUTHOR NAME(S) -- USER INPUT}\\
\small\texttt{AFFILIATION(S), CITY, COUNTRY -- USER INPUT}\\
\small\texttt{CORRESPONDING-AUTHOR EMAIL / ORCID -- USER INPUT}}
\date{September 2026}
\begin{document}
\maketitle

\begin{abstract}
For a nonconstant rational five-term arithmetic progression, we ask how many
nonzero terms can, after multiplication by one common rational scalar, be
cubes in a single nontrivial pure cubic field.  We prove that the exact
maximum is four.  The proof begins by determining the kernel of the map from
rational cube classes to the Kummer group of $\Q(\sqrt[3]{D})$.  After
normalizing radicands and quotienting by affine color change and reversal, we
obtain exactly 25 five-position color orbits.  Ten orbits are excluded by
known theorems on rational cubes in arithmetic progression.  A prime-support
lemma reduces the remaining 15 orbits to 60 explicit systems of three
diagonal cubic equations.  For every system, exhaustive finite-field
enumeration gives an obstruction at a displayed prime of good reduction; a
versioned certificate records all $p^2-1$ nonzero parameter pairs.  The
progression $(-3,-1,1,3,5)$ supplies four counted terms over
$\Q(\sqrt[3]{3})$.  Thus $R^{\times}_{(3,1)}(5)=4$.  This finite
classification excludes every five-hit color pattern but does not classify
all rational progressions attaining four hits.
\end{abstract}

\noindent\textbf{Keywords:} arithmetic progression; pure cubic field;
Kummer class; perfect cube; local obstruction.

\medskip
\noindent\textbf{Mathematics Subject Classification (2020):}
Primary 11B25; Secondary 11D41, 11R16, 11Y50.

\section{Definition and normalizations}
Let $A_i=a+id$ for $0\leq i<5$, where $a,d\in\Q$ and $d\ne0$.  For a
nontrivial class $[D]\in\Q^*/\Q^{*3}$ put
$K_D=\Q(\sqrt[3]{D})$.  Define $R^{\times}_{(3,1)}(5)$ as the maximum number
of nonzero $A_i$ for which there are one $D$ and one $\lambda\in\Q^*$ such
that $\lambda A_i\in K_D^3$ at every counted position.

The radicand may be chosen as a positive cube-free integer: $-1$ is already
a rational cube, and multiplying $D$ by a rational cube changes no field.
Moreover $D$ and $D^2$ generate the same field.  We quotient progressions by
common rational scaling and by reversal
$(a,d)\mapsto(a+4d,-d)$.  Zero has no class in $\Q^*/\Q^{*3}$ and is never
counted.  If $D$ is a rational cube, then $K_D=\Q$ and the problem is the
ordinary rational-cube problem, whose five-term maximum is three
\cite{HT}; this degree-one degeneration cannot attain our maximum.

\section{The pure-cubic Kummer kernel}
\begin{lemma}\label{lem:kernel}
If $[D]\ne1$ in $\Q^*/\Q^{*3}$, then
\[
\ker\bigl(\Q^*/\Q^{*3}\longrightarrow K_D^*/K_D^{*3}\bigr)
=\langle[D]\rangle=\{1,[D],[D]^2\}.
\]
Equivalently, a nonzero rational $r$ is a cube in $K_D$ if and only if
$[r]\in\langle[D]\rangle$.
\end{lemma}
\begin{proof}
Write $\alpha^3=D$ and $y=a+b\alpha+c\alpha^2$.  Reduction of $y^3$ in the
basis $1,\alpha,\alpha^2$ gives nonconstant coefficients
\begin{align*}
3(a^2b+Da c^2+Db^2c),\qquad
3(a^2c+ab^2+Dbc^2).
\end{align*}
If $a=0$, their vanishing implies $bc=0$, so $y$ lies on either the
$\Q\alpha$ or $\Q\alpha^2$ axis.  If $a\ne0$, put $B=b/a$, $C=c/a$.  The two
equations are
\[
B+DC^2+DB^2C=0,\qquad C+B^2+DBC^2=0.
\]
Their resultants, eliminating $C$ and $B$ respectively, are
\[
BD(B^3D-1)^2,\qquad C(C^3D^2-1)^2.
\]
Since $D$ is not a rational cube, the first forces $B=0$, and then the first
equation forces $C=0$.  Hence an element of $K_D$ with rational cube lies in
exactly one of $\Q$, $\Q\alpha$, $\Q\alpha^2$.  Cubing proves the asserted
three classes, and the converse is immediate.
\end{proof}

Thus, after the common scaling, every counted term has a unique color
$c_i\in\F_3$ and can be written
\begin{equation}\label{eq:colors}
A_i=D^{c_i}x_i^3,\qquad x_i\in\Q^*.
\end{equation}
Changing the common scale by $D^k$, replacing $D$ by $D^2$, and reversing
the progression act on color words by
\[
c_i\mapsto c_i+k,\qquad c_i\mapsto-c_i,\qquad
(c_0,\ldots,c_4)\mapsto(c_4,\ldots,c_0).
\]

\section{The 25 color orbits}
Exact enumeration gives 25 orbits of $\F_3^5$ under these operations.  For
completeness, Burnside's lemma gives the same count without trusting the
choice of representatives: the fixed-point counts of the 12 group elements
are
\[
243,27,9,9,9,1,1,1,0,0,0,0.
\]
Their sum is 300, and hence there are $300/12=25$ orbits.  Nine
contain a monochromatic three-term arithmetic progression of indices among
\[
(0,1,2),(1,2,3),(2,3,4),(0,2,4).
\]
Such an orbit would give three nonconstant rational cubes in arithmetic
progression after scaling, contrary to Darmon--Merel \cite{DM}.  Exactly one
further orbit, represented by $00100$, survives this test but has four terms
of one color, at positions $0,1,3,4$.  It contradicts the exact ordinary
bound $P_5(3)=3$ of Hajdu--Tengely \cite{HT}.  The remaining representatives
are listed below.  There is no hidden integrality assumption here: choosing
an integer $L$ that clears the finitely many rational cube roots and the
progression parameters, multiplication of the entire progression by the
common cube $L^3$ makes both the progression and those roots integral.
Cube positions and nonconstancy are therefore preserved.
\begin{gather*}
00101,00102,00110,00112,00120,00121,00122,01001,01002,01012,\\
01021,01102,01120,01201,01210.                         \tag{1}
\end{gather*}

\section{Prime support and the 60 curves}
Clear denominators and divide by the common gcd, obtaining a primitive
integer progression $A_i=a+id$.  Its classes have the affine form
$g+c_i\delta$ in $\Q^*/\Q^{*3}$.  For $i\ne j$,
\[
\gcd(A_i,A_j)\mid |i-j|,
\]
because $\gcd(a,d)=1$.  Therefore a prime $p>3$ divides at most one of the
five terms.

\begin{lemma}[prime-support lemma]\label{lem:support}
For each of the 15 words in (1), every $p>3$ occurs in every $A_i$ with
exponent divisible by three.  Consequently both $g$ and $\delta$ are
supported only at $2$ and $3$.
\end{lemma}
\begin{proof}
If a $p>3$ exponent were nonzero modulo three, the five valuation residues
would be supported at one position.  But they must be the values of the
affine function $c\mapsto g_p+c\delta_p$ on the displayed color word.
For a two-color word in (1), both color blocks have size at least two, so a
singleton support is impossible.  For a three-color word, an affine function
over $\F_3$ vanishing at the other two colors vanishes identically.  This is
again impossible.  Hence all five residues vanish.  Taking differences of
two used colors proves the assertion for $\delta$, and then for $g$.
\end{proof}

After another common rational scaling removes $g$, the four possible
one-dimensional subspaces of the $\{2,3\}$-supported class plane have
representatives
\[
D=2,3,6,18.
\]
For a word $c=(c_0,\ldots,c_4)$ and one such $D$, all five hits give the
concrete projective curve in $\mathbf P^4$
\begin{equation}\label{eq:curve}
D^{c_i}x_i^3-2D^{c_{i+1}}x_{i+1}^3
 +D^{c_{i+2}}x_{i+2}^3=0\quad(0\leq i\leq2).
\end{equation}

At every prime used below, $p\geq7$ and $p\nmid D$.  These conditions really
give good reduction.  Put $u_i=D^{c_i}x_i^2$.  After deleting the invertible
factor 3, the Jacobian matrix of the three equations is
\[
\begin{pmatrix}
u_0&-2u_1&u_2&0&0\\
0&u_1&-2u_2&u_3&0\\
0&0&u_2&-2u_3&u_4
\end{pmatrix}.
\]
A projective solution has at most one $u_i=0$: two zero terms in the induced
five-term progression would force both its initial term and common difference
to vanish because $p>4$.  If no $u_i$ is zero, columns $0,1,2$ give a nonzero
minor.  If the unique zero is respectively $u_0,u_1,u_2,u_3,u_4$, columns
$124,024,014,014,013$ give
\[
3u_1u_2u_4,-2u_0u_2u_4,u_0u_1u_4,u_0u_1u_4,-2u_0u_1u_3.
\]
Thus the Jacobian has rank 3 at every geometric point, proving smoothness.

Table~\ref{tab:local-obstructions} gives, in word order (1), a prime with no
nonzero projective point on \eqref{eq:curve}; columns correspond to
$D=2,3,6,18$.
\begin{table}[ht]
\centering\small
\begin{tabular}{c|rrrr@{\qquad}c|rrrr}
\toprule
word&2&3&6&18&word&2&3&6&18\\\midrule
00101&7&7&31&7&00102&7&7&13&7\\
00110&7&7&13&7&00112&7&7&13&7\\
00120&7&13&13&31&00121&7&19&61&31\\
00122&13&13&13&43&01001&7&7&13&7\\
01002&7&7&61&7&01012&13&7&13&7\\
01021&7&13&13&31&01102&13&7&13&7\\
01120&7&7&61&7&01201&7&7&13&7\\
01210&7&13&31&31&&&&&\\
\bottomrule
\end{tabular}
\caption{Good primes obstructing the 60 remaining word--radicand models}
\label{tab:local-obstructions}
\end{table}
Every listed prime is coprime to $3D$.  The certificate enumerates all
$(a,d)\in\F_p^2\setminus\{(0,0)\}$ and tests whether
$a+id\in D^{c_i}\F_p^3$ for all five positions.  This is equivalent to
enumerating the projective solutions of \eqref{eq:curve}, including solutions
with some zero coordinates.  Absence modulo $p$ excludes a rational point.

The complete finite computation is identified by
\begin{quote}\small\ttfamily
PAPER\_CUBE\_KUMMER5\_CERTIFICATE.json,\par
schema paper-cube-pure-cubic-kummer-n5-v2,\par
SHA-256\par
7c5ff0bd36ebfc3bace0ca5898625b2fd7f03d318c754c6650d3a2483dd83977.
\end{quote}
Its generator and tests are located by the root-level file
\begin{quote}\small\ttfamily
PAPER\_CUBE\_KUMMER5\_SUPPLEMENT\_MANIFEST.json.
\end{quote}
Each of the 60
cells scans all $p^2-1$ nonzero parameter pairs, for 23520 pairs in total.
The local supplement is not a public archive; at submission its manifest must
be deposited at the immutable DOI or URL stated in the data-availability
record.

\section{Exact maximum and scope of classification}
\begin{theorem}\label{thm:main}
With nonzero terms and nontrivial pure cubic fields as in Section 1,
\[
R^{\times}_{(3,1)}(5)=4.
\]
\end{theorem}
\begin{proof}
A five-hit progression has one of the 25 color orbits.  The first ten are
excluded in Section 3.  Lemma \ref{lem:support} and the 60 good-prime
obstructions exclude the remaining 15.  Hence the maximum is at most four.
For the reverse inequality, in $K=\Q(\alpha)$ with $\alpha^3=3$, the nonzero
progression
\[
-3,-1,1,3,5
\]
has its first four terms equal to the cubes of
$-\alpha,-1,1,\alpha$.
\end{proof}

The theorem includes a complete classification and exclusion of every
\emph{five-hit color/position class}.  It does not classify every rational
progression attaining four hits.  That latter set is a distinct Diophantine
problem: already a four-position pattern is naturally a weighted complete
intersection, and several patterns have multiple rational examples.  We do
not replace that rational-point classification by a bounded search.

If zero were counted, $-2,-1,0,1,2$ would give five hits in
$\Q(\sqrt[3]{2})$; this is exactly why the superscript $\times$ is necessary.

\section*{Statements and Declarations}

\paragraph{Funding.}
\texttt{FUNDING STATEMENT -- HUMAN AUTHOR INPUT REQUIRED.}

\paragraph{Competing interests.}
\texttt{COMPETING-INTERESTS STATEMENT -- HUMAN AUTHOR INPUT REQUIRED.}

\paragraph{Author contributions.}
\texttt{AUTHOR-CONTRIBUTION STATEMENT -- HUMAN AUTHOR INPUT REQUIRED.}

\paragraph{Data and code availability.}
The exact source, tests, and finite certificate will be deposited with the
accepted release at
\texttt{PERSISTENT ARCHIVE DOI/URL -- HUMAN AUTHOR INPUT REQUIRED}.  The local
release candidate is not yet a public archive.

\paragraph{Use of generative AI.}
Generative AI tools were used during exploratory mathematical research,
drafting, code generation, and preparation of the reproducibility package.
No AI system is listed as an author.  Before submission, the human authors
must independently verify every statement, proof, computation, citation, and
file, approve that this disclosure is complete, and accept full responsibility
for the submitted work.

\paragraph{Ethics approval, consent to participate, and consent for publication.}
Not applicable: this is a mathematical study involving no human participants,
animals, or personal data.

\paragraph{Supplementary information.}
Online Resource 1, \texttt{PAPER\_CUBE\_KUMMER5\_ARCHIVE\_CANDIDATE.zip}
(ZIP), contains the exact source, PDF, generator, tests, 60-model certificate,
manifest, data dictionary, and release metadata used for the finite
classification.  It is a local candidate until the human authors approve and
deposit the immutable archive.

\begin{thebibliography}{9}
\bibitem{DM} H.~Darmon and L.~Merel,
\emph{Winding quotient and some variants of Fermat's Last Theorem},
J. Reine Angew. Math. \textbf{490} (1997), 81--100,
\url{https://doi.org/10.1515/crll.1997.490.81}.
\bibitem{HT} L.~Hajdu and S.~Tengely,
\emph{Powers in arithmetic progressions}, Ramanujan J. \textbf{55}
(2021), 965--986, \url{https://doi.org/10.1007/s11139-020-00331-5}.
\end{thebibliography}
\end{document}
